% Greek abstract
\begin{abstract}
    Η πρόβλεψη ακμών σε γραφήματα με κειμενικά χαρακτηριστικά κόμβων αποτελεί το πρόβλημα
    εξαγωγής συμπεράσματος για την ύπαρξη ακμής μεταξύ δύο κόμβων, λαμβάνοντας υπόψη
    τόσο την τοπολογία του γραφήματος όσο και το κειμενικό περιεχόμενο των κόμβων.
    Πρόσφατες εργασίες σε τέτοια γραφήματα έχουν αναφέρει σχεδόν τέλεια ακρίβεια με
    καθαρά σημασιολογικές προσεγγίσεις βασισμένες σε προ-εκπαιδευμένα γλωσσικά
    μοντέλα, όμως το τετραγωνικό κόστος συμπερασμού τους τα καθιστά μη πρακτικά σε
    κλίμακα.

    Η παρούσα διπλωματική εργασία προτείνει το \textbf{\en{CascadeLP}}, ένα
    τετραεπίπεδο πλαίσιο πρόβλεψης ακμών που αντιμετωπίζει το κόστος συμπερασμού ως
    πρωταρχικό σχεδιαστικό περιορισμό: δρομολογεί κάθε ζεύγος κόμβων μέσω προοδευτικά
    πιο δαπανηρών μοντέλων βάσει ορίων αξιοπιστίας, κλιμακώνοντας μόνο όταν το
    φθηνότερο επίπεδο δεν είναι αρκετά σίγουρο.
    Η αλυσίδα περιλαμβάνει: αυτοβρόχους,
    δομικούς ευρετικούς δείκτες (\en{Common Neighbours}, \en{Jaccard},
    \en{Adamic-Adar}, \en{Preferential Attachment}), συχνότητες μερών του λόγου με
    \en{Random Forest} για κόμβους χωρίς δομικούς γείτονες, και έναν κωδικοποιητή
    \en{Sentence-Transformer} για ζεύγη στα οποία τα προηγούμενα επίπεδα εμφανίζουν
    αβεβαιότητα.

    Αξιολογούμε το \en{CascadeLP} στο σύνολο αξιολόγησης \en{DSAA 2023} της
    \en{Wikipedia}, το οποίο επιλέγεται ως πεδίο αξιολόγησης επειδή συνιστά
    ρεαλιστικό στιγμιότυπο γραφήματος με κειμενικά χαρακτηριστικά σε κλίμακα
    εκατοντάδων χιλιάδων κόμβων, με καθιερωμένο πίνακα κατάταξης διαγωνισμού που
    παρέχει υπαρκτές λύσεις προς σύγκριση.
    Για να επικυρώσουμε ότι η επιτυγχανόμενη
    ακρίβεια είναι γνήσια και όχι τεχνούργημα διαρροής δεδομένων ή τετριμμένης
    διαχωρισιμότητας, διεξάγουμε έναν συστηματικό έλεγχο \en{data leakage} --- που
    βρίσκει αμελητέα επικάλυψη ζευγών --- και έναν ανεξάρτητο χαρακτηρισμό
    διαχωρισιμότητας του συνόλου δεδομένων, τον οποίο διασταυρώνουμε ανά ζεύγος με το
    επίπεδο που το επέλυσε.
    Το \en{CascadeLP} πετυχαίνει \en{Macro F1} = \CascadeValFone{}
    σε πλήρως ανεξάρτητο σύνολο επικύρωσης --- επιβεβαιωμένο ανεξάρτητα με
    \en{Macro F1} = \KaggleBaselinePublic{} στο δημόσιο \en{leaderboard} του
    διαγωνισμού --- καλώντας τον ακριβό κωδικοποιητή \en{Sentence-Transformer} σε
    μόλις \TierThreePct\% των ζευγών.
    Συμπεριλαμβάνουμε επίσης ένα μοντέλο αναφοράς
    \en{TF-IDF + SVM}.
    Η γενίκευση σε δεύτερο σύνολο δεδομένων αναβάλλεται ως
    μελλοντική εργασία.

    \begin{keywords}
        Πρόβλεψη ακμών, γραφήματα με κειμενικά χαρακτηριστικά, νευρωνικά δίκτυα γραφημάτων,
        γλωσσικά μοντέλα, βαθμιδωτή συμπερασματολογία, \en{cold-start},
        \en{Wikipedia}, \en{DSAA 2023}
    \end{keywords}
\end{abstract}

% English abstract
\begin{abstracteng}
    \en{
        Link prediction in text-attributed graphs (TAGs) is the problem of inferring whether an
        edge exists between two nodes given both the graph topology and the textual content of
        the nodes. Recent work on such graphs has reported near-perfect accuracy with pure
        semantic approaches based on pre-trained language models, but their quadratic inference
        cost makes them impractical at scale.

        This thesis proposes \textbf{CascadeLP}, a four-tier link prediction framework that
        treats inference cost as a first-class design constraint: it routes each node pair
        through progressively more expensive models based on confidence thresholds, escalating
        only when the cheaper tier is not confident enough. The cascade includes: self-loops,
        structural heuristics (Common Neighbours, Jaccard, Adamic-Adar, Preferential
        Attachment), POS-tag frequency features with a Random Forest for nodes with no
        structural neighbours, and a Sentence-Transformer encoder for pairs where the earlier
        tiers remain uncertain.

        We evaluate CascadeLP on the DSAA 2023 Wikipedia benchmark, selected as our evaluation
        scope because it is a realistic, large-scale text-attributed graph with an established
        competition leaderboard of existing solutions to benchmark against. To validate that the
        achieved accuracy is genuine rather than an artifact of data leakage or trivial dataset
        separability, we run a systematic leakage audit -- which finds negligible pair overlap --
        and an independent dataset separability characterization, cross-referenced per pair with
        the tier that resolved it. CascadeLP achieves Macro F1 = \CascadeValFone{} on a fully
        held-out validation set -- independently confirmed with Macro F1 = \KaggleBaselinePublic{}
        on the competition's public leaderboard -- while calling the expensive
        Sentence-Transformer encoder on only \TierThreePct\% of pairs. We also include a TF-IDF
        + SVM reference baseline. Generalization to a second dataset is deferred to
        future work. }

    \begin{keywordseng}
        \en{Link prediction, text-attributed graphs, graph neural networks, language models,
            cascaded inference, cold-start, Wikipedia, DSAA 2023}
    \end{keywordseng}
\end{abstracteng}
